\clearpage\fancyhead[L]{\small ANZEIGE.}\pdfbookmark[0]{Anzeigen. 1815 Juli 1.}{anzeige-1815-pfaff}\SourceLeaf{240}
\begin{center}\textbf{A N Z E I G E N.}\par\bigskip
\rule{.85\linewidth}{.4pt}\par
Göttingische gelehrte Anzeigen.\quad1815 Juli 1.\par
\rule{.85\linewidth}{.4pt}\end{center}
Der Königl. Societät ist durch Hrn. Prof. \textsc{Gauss} eine handschriftliche Abhandlung des Hrn. Prof. \textsc{Pfaff} in Halle vorgelegt, überschrieben:
\begin{center}\textit{Methodus generalis, aequationes differentiarum partialium, nec non aequationes differentiales vulgares, utrasque primi ordinis, inter quotcunque variabiles, complete integrandi.}\end{center}
Die Lehre von den partiellen Differentialgleichungen des ersten Grades verdankt bekanntlich \textsc{Lagrange} zwei wichtige Erweiterungen, nemlich die allgemeine Integration derselben, wenn sie entweder, bei einer beliebigen Anzahl veränderlicher Grössen, die partiellen Differentialquotienten bloss linearisch enthalten, oder ohne Einschränkung in Rücksicht der Form, wenn der veränderlichen Grössen nur drei sind. Ueber diese beiden Fälle war man eigentlich bisher noch nicht hinausgegangen, und das von \textsc{Lagrange} bei den partiellen Differentialgleichungen dreier veränderlicher Grössen angewandte Verfahren würde bei einer grössern Anzahl besondere Schwierigkeiten haben. Wir haben es daher als eine merkwürdige Bereicherung der Integralrechnung anzusehen, dass es dem scharfsinnigen Verfasser der vorliegenden Abhandlung gelungen ist, die allgemeine Integration der partiellen Differentialgleichungen des ersten Grades für \textit{jede} Anzahl von veränderlichen Grössen zu finden. Er hat bei dieser Untersuchung einen ei\SourceLeaf{241}%
genthümlichen Weg gewählt, und sie an einen andern nicht weniger interessanten Zweig der Integralrechnung angeknüpft, nemlich an die Lehre von den gewöhnlichen Differentialgleichungen (des ersten Grades) zwischen mehr als zwei veränderlichen Grössen, deren wahre Natur bekanntlich erst \textsc{Monge} uns kennen gelehrt hat, obwohl die Integration derselben von diesem Geometer nur für die einfachsten Fälle vollendet ist. Von dieser Gattung von Differentialgleichungen giebt Hr. \textsc{Pfaff} die allgemeine Integration, und die der partiellen Differentialgleichungen erscheint dann nur als ein besonderer Fall von jener. Bei diesen Untersuchungen wird die allgemeine Integration der Differentialgleichungen von jedem Grade zwischen zwei veränderlichen Grössen \textit{vorausgesetzt}, welches ganz in der Ordnung ist, eben so wie man in den höhern Theilen der Mathematik die allgemeine Auflösung der algebraischen Gleichungen postulirt.

Wir glauben den Freunden der höhern Mathematik einen angenehmen Dienst zu erweisen, wenn wir sie durch gegenwärtige Anzeige in den Besitz dieser schönen Erweiterung der Integralrechnung setzen. Freilich würde ein Auszug aus der 144 Quartseiten starken Abhandlung, in welchem wir dem Verfasser Schritt für Schritt folgen und nichts Wesentliches übergehen wollten, die Grenzen des uns vergönnten Raumes weit überschreiten. Wir wollen daher versuchen, indem wir uns bloss an die \textit{Sache} halten, in einer etwas veränderten Darstellung das Wesentliche so herauszuheben, dass Kenner sich dasselbe vollkommen aneignen können.

Als die Hauptoperation des ganzen Geschäfts muss angesehen werden die Reduction eines Differentialausdrucks
\[
\Omega=p\,\mathrm dx+p'\mathrm dx'+p''\mathrm dx''+\text{etc.}+p^{(n-1)}\mathrm dx^{(n-1)}
\]
wo jeder der Coëfficienten $p,p',p''$ u. s. w. Function der veränderlichen Grössen $x,x',x''$ u. s. w. ist auf die Form
\[
\Omega=\lambda(q\,\mathrm dy+q'\mathrm dy'+q''\mathrm dy''+\text{etc.}+q^{(n-2)}\mathrm dy^{(n-2)});
\]
so dass $\lambda,y,y',y''$ u. s. f. Functionen von $x,x',x''$ u. s. w. seien, hingegen $q,q',q''$ u. s. f. Functionen von $y,y',y''$ u. s. w., und dass die Anzahl der letztern veränderlichen Grössen um eine kleiner sei, als die Anzahl der veränderlichen Grössen $x,x',x''$ u. s. f.

Diese Verwandlung, welche wir Kürze halber mit (I) bezeichnen wollen, beschränkt sich auf den Fall, wo $n$ eine gerade Zahl ist; wir werden weiter un\SourceLeaf{242}%
ten entwickeln, wie sie ausgeführt werden müsse, und sie, um die Uebersicht
nicht zu stören, hier einstweilen voraussetzen.
\par
In dem Fall. wo $n$ eine ungerade Zahl ist, würde die Verwandlung I nur
unter speciellen Bedingungen zwischen den Coëfficienten $p,p',p''.\,.\,.$ möglich
sein; allgemein aber lässt sich in diesem Falle $\Omega$ auf die Form
\[
p^{*}\mathrm{d}x+\lambda(q\mathrm{d}y+q'\mathrm{d}y'+q''\mathrm{d}y''+\mathrm{etc.}
+q^{(n-3)}\mathrm{d}y^{(n-3)})
\]
bringen. Man sehe nemlich einstweilen $x$ in $\Omega$ als constant an, und verwandle
unter dieser Voraussetzung nach (I)
\[
p'\mathrm{d}x'+p''\mathrm{d}x''+p'''\mathrm{d}x'''+\mathrm{etc.}
+p^{(n-1)}\mathrm{d}x^{(n-1)}
\]
wo nunmehr die Anzahl der veränderlichen Grössen $x',x'',x'''.\,.\,.$ gerade sein
wird, in
\[
\lambda(q\mathrm{d}y+q'\mathrm{d}y'+q''\mathrm{d}y''+\mathrm{etc.}
+q^{(n-3)}\mathrm{d}y^{(n-3)})=\lambda\Omega'
\]
Hier werden also $q,q',q''$ u.\,s.\,w. Functionen von $y,y',y''$ u.\,s.\,w. sein, diese hingegen, eben so wie $\lambda$, Functionen von $x,x',x'',x'''.\,.\,.\,.$, von welchen Grössen
jedoch die erste $x$ als constant behandelt werden muss, um aus der Entwickelung von $\lambda\Omega'$
\[
p'\mathrm{d}x'+p''\mathrm{d}x''+p'''\mathrm{d}x'''+\mathrm{etc.}
\]
zu erhalten. Das Glied, welches noch hinzukommt, wenn bei jener Entwicklung
auch $x$ als veränderlich betrachtet wird, ist
\[
=(q\cdot\frac{\mathrm{d}y}{\mathrm{d}x}
+q'\cdot\frac{\mathrm{d}y'}{\mathrm{d}x}
+q''\cdot\frac{\mathrm{d}y''}{\mathrm{d}x}+\mathrm{etc.})\lambda\mathrm{d}x
\]
Man hat daher, um die obige Form zu erhalten, nur
\[
p^{*}=p-\lambda(q\cdot\frac{\mathrm{d}y}{\mathrm{d}x}
+q'\cdot\frac{\mathrm{d}y'}{\mathrm{d}x}
+q''\cdot\frac{\mathrm{d}y''}{\mathrm{d}x}+\mathrm{etc.})
\]
zu setzen. Diese Verwandlung von $\Omega$ in $p^{*}\mathrm{d}x+\lambda\Omega'$, welche auf ungerade Werthe von $n$ beschränkt ist, wollen wir mit II bezeichnen. Offenbar kann dieselbe
Reduction abermals auf $\Omega'$ angewandt und
\[
\begin{aligned}
&\Omega'=q^{*}\mathrm{d}y+\lambda'(r\mathrm{d}z+r'\mathrm{d}z'+r''\mathrm{d}z''+\mathrm{etc.}\\&\qquad+r^{(n-5)}\mathrm{d}z^{(n-5)})\\
&=q^{*}\mathrm{d}y+\lambda'\Omega''
\end{aligned}
\]
\SourceLeaf{243}
gesetzt werden, und so abermals $\Omega''=r^{*}\mathrm{d}z+\lambda''\Omega'''$ bis man zuletzt auf einen
Ausdruck kommt, der bloss Eine veränderliche Grösse enthält. Dadurch ist also
$\Omega$ auf die Form
\[
p^{*}\mathrm{d}x+\lambda q^{*}\mathrm{d}y+\lambda\lambda'r^{*}\mathrm{d}z+\mathrm{etc.}
\]
gebracht, oder auf die Form
\[
P\mathrm{d}x+Q\mathrm{d}y+R\mathrm{d}z+\mathrm{etc.}
\]
wo die Anzahl der veränderlichen Grössen $x,y,z$ u.\,s.\,w. $=\frac12(n+1)$, und wo die
sämmtlichen $n$ Grössen $y,z$ u.\,s.\,w. $P,Q,R$ u.\,s.\,w. Functionen von $x,x',x''$ u.\,s.\,w.
sein werden. Diess Reductionsverfahren mag durch III bezeichnet werden.

Wendet man diess Verfahren III in dem Fall, wo ursprünglich eine gerade
Anzahl veränderlicher Grössen vorgegeben war, auf den durch die Reduction I
erhaltenen Ausdruck
\[
q\mathrm{d}y+q'\mathrm{d}y'+q''\mathrm{d}y''+\mathrm{etc.}
\]
an, so kommt dadurch
\[
\begin{aligned}&\Omega=p\mathrm{d}x+p'\mathrm{d}x'+p''\mathrm{d}x''+\mathrm{etc.}
+p^{(n-1)}\mathrm{d}x^{(n-1)}\\&\qquad\text{in die Form}\end{aligned}
\]
\[
Q\mathrm{d}y+R\mathrm{d}z+\mathrm{etc.}
\]
so dass die Anzahl der veränderlichen Grössen $y,z$ u.\,s.\,w. $=\frac12 n$ wird. und alle
$n$ Grössen $Q,R$ u.\,s.\,f. $y,z$ u.\,s.\,f. Functionen von $x,x',x''$ u.\,s.\,w. werden.
Diese Reduction werde mit IV bezeichnet.

Diese allgemeine Transformabilität der Differentialausdrücke nach III und
IV ist ein eben so neuer als merkwürdiger Lehrsatz, der sich zwar in der Abhandlung des Hrn \textsc{Pfaff} nicht ausdrücklich ausgesprochen findet, aber sich leicht
aus den dortigen Untersuchungen folgern lässt.

Es lassen sich nun daraus die Auflösungen der im Eingange dieser Anzeige
erwähnten Aufgaben mit Leichtigkeit ableiten.

\medskip
1) Um die Differentialgleichung
\[
0=p\mathrm{d}x+p'\mathrm{d}x'+p''\mathrm{d}x''+\mathrm{etc.}
\]
oder
\[
0=\Omega
\]
\SourceLeaf{244}
zu integriren, wo $p,p',p''$ u.\,s.\,w. gegebne Functionen der $n$ veränderlichen
Grössen $x,x',x''$ u.\,s.\,w. sind, wird man, wenn $n$ gerade ist nach IV
\[
\Omega=Q\mathrm{d}y+R\mathrm{d}z+S\mathrm{d}u+\mathrm{etc.}
\]
machen, wo $Q,y.R,z,S,u$ u.\,s.\,w. zusammen $n$ gegebne Functionen von $x,x',x''$
sein werden. Die Differentialgleichung
\[
0=Q\mathrm{d}y+R\mathrm{d}z+S\mathrm{d}u+\mathrm{etc.}
\]
wird also der vorgegebnen gleichgeltend und ihre allgemeinste Integration in folgendem System von $\frac12 n$ Gleichungen enthalten sein:
\[
\begin{aligned}&0=\varphi(y,z,u\ \mathrm{etc.}),\\&
\frac1Q\cdot\frac{\mathrm{d}\varphi(y,z,u.\,.\,.)}{\mathrm{d}y}
\\&=\frac1R\cdot\frac{\mathrm{d}\varphi(y,z,u.\,.\,.)}{\mathrm{d}z}
\\&=\frac1S\cdot\frac{\mathrm{d}\varphi(y,z,u.\,.\,.)}{\mathrm{d}u}=\mathrm{etc.}\end{aligned}
\]
wo $\varphi$ eine willkürliche Function vorstellt. und die Differentialquotienten, wie
sich von selbst versteht, \textit{partielle} sind. In so fern vermittelst der Gleichung
$0=\varphi(y,z,u\ \mathrm{etc.})$ die Grösse $y$ sich durch die übrigen bestimmen lässt, kann
man die Auflösung auch durch folgende Gleichungen darstellen:
\[
\begin{gathered}
y=\psi(z,u.\,.\,.\,.)\\
\frac{\mathrm{d}\psi(z,u.\,.\,.\,.)}{\mathrm{d}z}=-\frac RQ\\
\frac{\mathrm{d}\psi(z,u.\,.\,.\,.)}{\mathrm{d}u}=-\frac SQ\\
\text{u.\,s.\,f.}
\end{gathered}
\]
Genau genommen wäre indessen diese Auflösung weniger allgemein, da die willkürliche Function $\varphi(y,z,u.\,.\,.\,.)$ auch solche unter sich begreift, in welchen $y$
nicht mit vorkommt.

\medskip
2) Zur Integration derselben Differentialgleichung in dem Falle, wo $n$ ungerade ist, wird man $\Omega$ nach III in folgende Form setzen
\[
\Omega=P\mathrm{d}x+Q\mathrm{d}y+R\mathrm{d}z+\mathrm{etc.}
\]
wo $P,Q,y,R,z$ u.\,s.\,w. zusammen $n$ gegebne Functionen von $x,x',x''$ u.\,s.\,w.
sein werden. Die allgemeinste Integration der Differentialgleichung $\Omega=0$ beruhet dann auf folgendem System von $\frac12(n+1)$ Gleichungen:
\[
\begin{aligned}&0=\varphi(x,y,z.\,.\,.),\\&
\frac1P\cdot\frac{\mathrm{d}\varphi(x,y,z.\,.\,.)}{\mathrm{d}x}
=\frac1Q\cdot\frac{\mathrm{d}\varphi(x,y,z.\,.\,.)}{\mathrm{d}y}
\\&=\frac1R\cdot\frac{\mathrm{d}\varphi(x,y,z.\,.\,.)}{\mathrm{d}z}=\mathrm{etc.}\end{aligned}
\]
\par\SourceLeaf{245}
3) Die allgemeine Integration einer gegebenen partiellen Differentialgleichung
des ersten Grades, d.\,i. einer endlichen Gleichung zwischen den partiellen Differentialquotienten
\[
\frac{\mathrm{d}x}{\mathrm{d}x'}=p',\qquad
\frac{\mathrm{d}x}{\mathrm{d}x''}=p'',\qquad
\frac{\mathrm{d}x}{\mathrm{d}x'''}=p''',\quad\text{u.\,s.\,f.}
\]
und $x,x',x'',x'''$ u.\,s.\,w. (wo $x$ eine erst zu bestimmende Function der $m$ veränderlichen Grössen $x',x'',x'''$ u.\,s.\,w. vorstellt) ist nichts anders, als die allgemeine Integration der gewöhnlichen Differentialgleichung
\[
0=-\mathrm{d}x+p'\mathrm{d}x'+p''\mathrm{d}x''+p'''\mathrm{d}x'''+\text{etc.}
\]
Da nemlich vermöge jener endlichen Gleichung eine der Grössen $p',p'',p'''$ u.\,s.\,w.
z.\,B. $p'$ als Function der übrigen $p'',p'''$ u.\,s.\,w. und $x,x',x'',x'''$ u.\,s.\,w. dargestellt werden kann, so ist die eben angegebne Differentialgleichung als eine zwischen den $2m$ veränderlichen Grössen $x,x',x''$ u.\,s.\,w. $p'',p'''$ u.\,s.\,w. zu betrachten, in welcher die Differentiale $\mathrm{d}p'',\mathrm{d}p'''$ u.\,s.\,w. mit dem Coëfficienten $0$ behaftet sind. Um also die Integration auszuführen, wird man den Differentialausdruck
\[
-\mathrm{d}x+p'\mathrm{d}x'+p''\mathrm{d}x''+p'''\mathrm{d}x'''+\text{etc.}
\]
auf die Form
\[
Q\mathrm{d}y+R\mathrm{d}z+S\mathrm{d}u+\text{etc.}
\]
bringen, wo die $2m$ Grössen $Q,y,R,z,S.\ u$ u.\,s.\,w. bekannte Functionen von
$x,x',x''\, .\, .\, .\, .\, p'',p'''\, .\, .\, .\, .$ sein werden. Die Integration ist sodann in demselben
System von Gleichungen wie oben (1) enthalten, und wenn man sich aus ihnen
$p'',p'''$ u.\,s.\,w. eliminirt denkt, bleibt Eine endliche Gleichung zwischen $x,x',x''.x'''$
etc. zurück. Die \emph{wirkliche} Elimination kann freilich nur ausgeführt werden, in so
fern für $\varphi$ \emph{bestimmte} Functionen angenommen werden; allein dieser Umstand beruhet auf der Natur des Problems und nicht auf der Unvollkommenheit der Analyse, welche, so lange sie beim Allgemeinen stehen bleibt, die Auflösung nur in
jener Form geben kann.

Uebrigens sieht man von selbst, dass auf ähnliche Art die Integration \emph{mehrerer} neben einander bestehender partieller Differentialgleichungen in unsrer Gewalt ist.

Es bleibt uns jetzt nichts weiter übrig, als nur noch eine allgemeine Me\SourceLeaf{246}%
thode für die oben mit (I) bezeichnete Transformation anzugeben. Was für
Functionen von $x,x',x''$ u.\,s.\,w. auch immer für $y,y',y''$ u.\,s.\,w. angenommen
werden, so ist klar, dass, wenigstens allgemein zu reden, durch Elimination die
Grössen $x',x'',x'''$ u.\,s.\,w. sich als Functionen von $x,y,y',y''$ u.\,s.\,w. werden darstellen lassen, deren Differentiation $n-1$ Gleichungen hervorbringen wird:
\[
\begin{aligned}
\mathrm{d}x'&=\xi'\mathrm{d}x+\alpha'\mathrm{d}y+\beta'\mathrm{d}y'+\gamma'\mathrm{d}y''+\text{etc.}\\
\mathrm{d}x''&=\xi''\mathrm{d}x+\alpha''\mathrm{d}y+\beta''\mathrm{d}y'+\gamma''\mathrm{d}y''+\text{etc.}\\
\mathrm{d}x'''&=\xi'''\mathrm{d}x+\alpha'''\mathrm{d}y+\beta'''\mathrm{d}y'+\gamma'''\mathrm{d}y''+\text{etc.}\quad\text{u.\,s.\,w.}
\end{aligned}
\]
Hier sind also die Coëfficienten $\xi',\xi'',\xi'''$ u.\,s.\,w. $\alpha',\alpha'',\alpha'''$ u.\,s.\,w. Functionen von
$x,y,y',y''$ u.\,s.\,w., und in dieser Beziehung werden wir ihre partiellen Differentialquotienten nach $x$ durch Einschliessung in Klammern unterscheiden: offenbar können jene Grössen auch als Functionen von $x,x',x'',x'''$ u.\,s.\,w. angesehen
werden, in welcher Beziehung wir den partiellen Differentialquotienten nach $x$
ohne Klammer schreiben wollen, so dass $\left(\frac{\mathrm{d}\xi'}{\mathrm{d}x}\right)$ wohl von $\frac{\mathrm{d}\xi'}{\mathrm{d}x}$ unterschieden werden muss. Dasselbe gilt von $\left(\frac{\mathrm{d}p}{\mathrm{d}x}\right)$ und $\frac{\mathrm{d}p}{\mathrm{d}x}$ u.\,s.\,w. Damit nun $\Omega$ nach Substitution jener Werthe von $\mathrm{d}x',\mathrm{d}x'',\mathrm{d}x'''$ u.\,s.\,w. die vorgeschriebene Form erhalte,
muss offenbar \emph{erstlich} $\mathrm{d}x$ herausfallen, also folgende Bedingungsgleichung [1]
Statt finden:
\[
0=p+p'\xi'+p''\xi''+p'''\xi'''+\text{etc.}
\]
Ferner sollen die Coëfficienten von $\mathrm{d}y,\mathrm{d}y',\mathrm{d}y''$ u.\,s.\,w. nemlich
\[
\begin{aligned}
p'\alpha'+p''\alpha''+p'''\alpha'''+\text{etc.}&=A\\
p'\beta'+p''\beta''+p'''\beta'''+\text{etc.}&=B\\
p'\gamma'+p''\gamma''+p'''\gamma'''+\text{etc.}&=C,\quad\text{u.\,s.\,w.}
\end{aligned}
\]
die Form $\lambda q,\lambda q',\lambda q''$ u.\,s.\,w. erhalten, so dass $q,q',q''$ u.\,s.\,w. bloss Functionen
von $y,y',y''$ u.\,s.\,w. werden; damit diess geschehe, müssen wir \emph{zweitens} haben [2]:
\[
\frac{1}{A}\cdot\left(\frac{\mathrm{d}A}{\mathrm{d}x}\right)
=\frac{1}{B}\cdot\left(\frac{\mathrm{d}B}{\mathrm{d}x}\right)
=\frac{1}{C}\cdot\left(\frac{\mathrm{d}C}{\mathrm{d}x}\right)\text{etc.}
=\frac{1}{\lambda}\cdot\left(\frac{\mathrm{d}\lambda}{\mathrm{d}x}\right)
\]
Nun ist aber
\[
\begin{aligned}
\left(\frac{\mathrm{d}A}{\mathrm{d}x}\right)={}&
p'\left(\frac{\mathrm{d}\alpha'}{\mathrm{d}x}\right)
+p''\left(\frac{\mathrm{d}\alpha''}{\mathrm{d}x}\right)
+p'''\left(\frac{\mathrm{d}\alpha'''}{\mathrm{d}x}\right)+\text{etc.}\\
&+\alpha'\left(\frac{\mathrm{d}p'}{\mathrm{d}x}\right)
+\alpha''\left(\frac{\mathrm{d}p''}{\mathrm{d}x}\right)
+\alpha'''\left(\frac{\mathrm{d}p'''}{\mathrm{d}x}\right)+\text{etc.}
\end{aligned}
\]
\SourceLeaf{247}
Substituirt man hier
\[
\left(\frac{\mathrm{d}\alpha'}{\mathrm{d}x}\right)=\frac{\mathrm{d}\xi'}{\mathrm{d}y},\qquad
\left(\frac{\mathrm{d}\alpha''}{\mathrm{d}x}\right)=\frac{\mathrm{d}\xi''}{\mathrm{d}y}.\quad\text{etc.}
\]
und subtrahirt
\[
\begin{aligned}
0={}&\frac{\mathrm{d}p}{\mathrm{d}y}
+\xi'\cdot\frac{\mathrm{d}p'}{\mathrm{d}y}
+\xi''\cdot\frac{\mathrm{d}p''}{\mathrm{d}y}+\text{etc.}\\
&+p'\cdot\frac{\mathrm{d}\xi'}{\mathrm{d}y}
+p''\cdot\frac{\mathrm{d}\xi''}{\mathrm{d}y}+\text{etc.}
\end{aligned}
\]
welche Gleichung entsteht, wenn man [1] nach $y$ differentiirt, so wird
\[
\left(\frac{\mathrm{d}A}{\mathrm{d}x}\right)=
\left\{\begin{aligned}
&\phantom{-\frac{\mathrm{d}p}{\mathrm{d}y}}+\alpha'\cdot\left(\frac{\mathrm{d}p'}{\mathrm{d}x}\right)
+\alpha''\cdot\left(\frac{\mathrm{d}p''}{\mathrm{d}x}\right)+\text{etc.}\\
&-\frac{\mathrm{d}p}{\mathrm{d}y}-\xi'\cdot\frac{\mathrm{d}p'}{\mathrm{d}y}
-\xi''\cdot\frac{\mathrm{d}p''}{\mathrm{d}y}-\text{etc.}
\end{aligned}\right.
\]
Da man nun ferner hat
\[
\begin{aligned}
\left(\frac{\mathrm{d}p'}{\mathrm{d}x}\right)&=\frac{\mathrm{d}p'}{\mathrm{d}x}
+\xi'\cdot\frac{\mathrm{d}p'}{\mathrm{d}x'}
+\xi''\cdot\frac{\mathrm{d}p'}{\mathrm{d}x''}+\text{etc.}\\
\left(\frac{\mathrm{d}p''}{\mathrm{d}x}\right)&=\frac{\mathrm{d}p''}{\mathrm{d}x}
+\xi'\cdot\frac{\mathrm{d}p''}{\mathrm{d}x'}
+\xi''\cdot\frac{\mathrm{d}p''}{\mathrm{d}x''}+\text{etc.},\quad\text{u.\,s.\,w.}
\end{aligned}
\]
\[
\begin{aligned}
\frac{\mathrm{d}p}{\mathrm{d}y}&=\alpha'\cdot\frac{\mathrm{d}p}{\mathrm{d}x'}
+\alpha''\cdot\frac{\mathrm{d}p}{\mathrm{d}x''}
+\alpha'''\cdot\frac{\mathrm{d}p}{\mathrm{d}x'''}+\text{etc.}\\
\frac{\mathrm{d}p'}{\mathrm{d}y}&=\alpha'\cdot\frac{\mathrm{d}p'}{\mathrm{d}x'}
+\alpha''\cdot\frac{\mathrm{d}p'}{\mathrm{d}x''}
+\alpha'''\cdot\frac{\mathrm{d}p'}{\mathrm{d}x'''}+\text{etc.},\quad\text{u.\,s.\,w.}
\end{aligned}
\]
so wird nach diesen Substitutionen
\[
\left(\frac{\mathrm{d}A}{\mathrm{d}x}\right)=k'\alpha'+k''\alpha''+k'''\alpha'''+\text{etc.}
\]
werden, wo
\[
\begin{array}{r@{}l@{}l@{}l@{}l@{}l}
k'={}&(1,0)&\quad *\quad&+(1,2)\xi''&+(1,3)\xi'''&+\text{etc.}\\
k''={}&(2,0)&+(2,1)\xi'&\quad *\quad&+(2,3)\xi'''&+\text{etc.}\\
k'''={}&(3,0)&+(3,1)\xi'&+(3.2)\xi''-&\quad *\quad&+\text{etc.}
\end{array}
\]
u.\,s.\,w. wenn man Kürze halber allgemein
\[
\frac{\mathrm{d}p^{(\mu)}}{\mathrm{d}x^{(\nu)}}
-\frac{\mathrm{d}p^{(\nu)}}{\mathrm{d}x^{(\mu)}}
\]
durch $(\mu,\nu)$ bezeichnet, so dass allgemein $(\mu,\mu)=0$, und $(\nu,\mu)=-(\mu.\nu)$ wird.
Ferner sieht man leicht, dass auch
\[
\begin{aligned}
\left(\frac{\mathrm{d}B}{\mathrm{d}x}\right)&=k'\beta'+k''\beta''+k'''\beta'''+\text{etc.}\\
\left(\frac{\mathrm{d}C}{\mathrm{d}x}\right)&=k'\gamma'+k''\gamma''+k'''\gamma'''+\text{etc.}
\end{aligned}
\]
\SourceLeaf{248}
u.\ s.\ w. wird, und dass folglich den Gleichungen [2] werde Genüge geleistet werden, wenn \(k',\,k'',\,k'''\) u.\ s.\ w. resp. den Grössen \(p'.\,p'',\,p'''\) u.\ s.\ w. proportional werden. Setzt man übrigens noch
\[
k=\quad *\quad +(0,1)\xi'+(0,2)\xi''+(0,3)\xi'''+\text{etc.}
\]
so hat man die identische Gleichung
\[
0=k+k'\xi'+k''\xi''+k'''\xi'''+\text{etc.}
\]
aus welcher mit [1] verbunden leicht gefolgert wird, dass auch \(k\) der Grösse \(p\) proportional sein muss; diese letztere Proportionalität kann die Stelle der Gleichung [1] vertreten. Mit Hülfe der \(n-1\) Gleichungen
\[
\frac{k}{p}=\frac{k'}{p'}=\frac{k''}{p''}=\frac{k'''}{p'''}=\text{u.\ s.\ w.}
\]
können nun die bisher unbekannten Functionen \(\xi',\,\xi'',\,\xi'''\) u.\ s.\ w., deren Anzahl gleichfalls \(n-1\) ist, bestimmt werden; jedoch zeigt eine nähere Betrachtung, dass diese Bestimmung nur für gerade Werthe von \(n\) ausführbar ist; für ungerade \(n\) wird allemal, sobald die Grössen \(\xi'.\,\xi'',\,\xi'''\) u.\ s.\ w. bis auf eine eliminirt sind, diese von selbst herausfallen und bloss eine \emph{Bedingungsgleichung} zwischen \(p,\,p',\,p''\) u.\ s.\ w. übrig bleiben. In diesem Umstande liegt der Grund, warum die Verwandlung (I) auf gerade Werthe von \(n\) beschränkt werden muss.

Die Bedingungen der Verwandlung I sind also jetzt darauf zurückgeführt, dass die partiellen Differentialquotienten \(\bigl(\frac{\mathrm{d}x'}{\mathrm{d}x}\bigr),\ \bigl(\frac{\mathrm{d}x''}{\mathrm{d}x}\bigr),\ \bigl(\frac{\mathrm{d}x'''}{\mathrm{d}x}\bigr)\) u.\ s.\ w., in so fern \(x',\,x'',\,x'''\) u\ s.\ w. als Functionen von \(x,\,y,\,y',\,y''\) u.\ s.\ w. betrachtet werden, die jetzt als bekannte Functionen von \(x,\,x',\,x''\) u.\ s.\ w. dargestellten Werthe \(\xi',\xi'',\xi'''\) u.\ s.\ w. erhalten. Diess lässt sich auch so ausdrücken: In so fern \(y,\,y',\,y''\) u.\ s.\ w. als constant und also \(x',\,x'',\,x'''\) u.\ s.\ w. bloss als Functionen der veränderlichen Grösse \(x\) betrachtet werden, muss folgenden \(n-1\) Differentialgleichungen Genüge geleistet werden
\[
\mathrm{d}x'=\xi'\mathrm{d}x,\qquad
\mathrm{d}x''=\xi''\mathrm{d}x,\qquad
\mathrm{d}x'''=\xi'''\mathrm{d}x,\quad\text{u.\ s.\ w.}
\]
oder wenn man die ursprünglichen Gleichungen vorzieht, aus deren Combination diese eigentlich entstanden waren, folgenden:
\SourceLeaf{249}
\[
\begin{aligned}
&\quad *\qquad +\frac{(0,1)}{p}\mathbin{.}\mathrm{d}x'
 +\frac{(0,2)}{p}\mathbin{.}\mathrm{d}x''
 +\frac{(0,3)}{p}\mathbin{.}\mathrm{d}x'''+\text{etc.}\\
&=\frac{(1,0)}{p'}\mathbin{.}\mathrm{d}x\qquad *\quad
 +\frac{(1,2)}{p'}\mathbin{.}\mathrm{d}x''
 +\frac{(1,3)}{p'}\,\mathrm{d}x'''+\text{etc.}\\
&=\frac{(2,0)}{p''}\mathbin{.}\mathrm{d}x
 +\frac{(2,1)}{p''}\mathbin{.}\mathrm{d}x'\qquad *\quad
 +\frac{(2,3)}{p''}\mathbin{.}\mathrm{d}x'''+\text{etc.}\\
&=\text{u.\ s.\ w.}
\end{aligned}
\]
Die Integration dieser Gleichungen gehört aber in das Gebiet der gewöhnlichen Integralrechnung, und wird hier vorausgesetzt; sie wird, allgemein zu reden, \(n-1\) von einander unabhängige Constanten enthalten \(H,\,H',\,H'',\,H'''\) u.\ s.\ w., die als gegebene Functionen von \(x',\,x'',\,x'''\) u.\ s.\ w. erscheinen, so dass man \(n-1\) endliche Gleichungen erhält
\[
H=X,\qquad H'=X',\qquad H''=X'',\quad\text{u.\ s.\ w.}
\]
wenn \(X,\,X',\,X''\) u.\ s.\ w. diese Functionen vorstellen. Es erhellt also aus dieser Analyse, dass den vorgeschriebenen Bedingungen Genüge geleistet sein wird, wenn man eben diese Functionen für \(y,\,y',\,y''\) u.\ s.\ w. wählt, oder
\[
y=X,\qquad y'=X',\qquad y''=X'',\quad\text{u.\ s.\ w.}
\]
setzt.

Eine Bemerkung wollen wir hier noch beifügen. Wir haben mit Vorbedacht gesetzt, dass die Integration, \emph{allgemein zu reden}, \(n-1\) von einander unabhängige Constanten gebe. In speciellen Fällen nemlich, d.\ i. wenn die Coefficienten \(p,\,p',\,p''\) u.\ s.\ w. so beschaffen sind, dass obige \(n-1\) Differentialgleichungen nicht von einander unabhängig sind, sondern eine schon aus Combination der übrigen abgeleitet werden kann, gilt diess nicht mehr: hier wird auch die Bestimmung von \(\xi',\,\xi'',\,\xi'''\) u.\ s.\ w. durch Elimination nicht mehr ausführbar sein. Dieser Fall müsste eigentlich als Ausnahme besonders behandelt werden; wir begnügen uns indessen hier um so mehr mit einer \emph{kurzen Andeutung}, wie man sich dabei verhalten könne, da der Verf. ihn nicht berührt hat. Man braucht nemlich an die Stelle der einen von obigen Differentialgleichungen, die schon in den andern enthalten ist, nur irgend eine andere willkürliche in diesen noch nicht enthaltene linearische Gleichung zwischen \(\mathrm{d}x,\,\mathrm{d}x',\,\mathrm{d}x''\) u.\ s.\ w. zu setzen, um das vorige anwenden zu können. Am bequemsten wird es immer sein, eines von diesen Differentialien \(=0\) zu setzen, oder eine von den veränderlichen Grössen \(x,\,x',\,x''\) u.\ s.\ w.
\SourceLeaf{250}
als constant zu behandeln, z. B. \(x\), wenn nicht zufällig \(\mathrm{d}x=0\) schon aus den
vorhandenen Gleichungen folgt. In diesem Fall wird eine von den Integralgleichungen sein
\[
H=x
\]
und wenn man \(y=x\) setzt, so lässt sich zeigen, dass in dem verwandelten Ausdrucke von \(\Omega\) allemal von selbst \(q=0\) wird.

\begin{center}\rule{0.25\linewidth}{0.3pt}\end{center}

